\documentclass{article}

\textwidth=6.5in
\textheight=8.5in
\voffset=-54pt
\oddsidemargin=0in
\evensidemargin=0in
\setlength{\parskip}{8pt}
\setlength{\parindent}{0pt}

\usepackage{amsmath, amssymb}
\usepackage{ifthen}

\usepackage[inline]{enumitem}
\setlist{topsep=0pt, parsep=8pt}

\usepackage[rgb,table]{xcolor}\convertcolorsUfalse
\usepackage{graphicx}

% change hyperref options
\PassOptionsToPackage{hyphens}{url}
\usepackage[bookmarks,colorlinks,linkcolor=black,citecolor=blue,pdfstartview=FitH,urlcolor=blue]{hyperref}
\hypersetup{pdfpagemode=UseNone}
\usepackage{bookmark}

\usepackage{graphicx}
\usepackage{multicol}
\usepackage{framed}
\usepackage{array}

\usepackage{icomma}

\usepackage{soul}
\usepackage[normalem]{ulem}

\usepackage{pgfplots}
\pgfplotsset{compat=1.18}  
\pgfplotscreateplotcyclelist{mycolor}{black, blue, red, green!70!black, magenta, purple, cyan, orange }  
\pgfplotsset{
  disabledatascaling,
  cycle list=tikzcolor, 
  axis lines=middle,
  every axis plot/.style={no marks, thick},
  every axis/.style={
    enlarge x limits=false,
    enlarge y limits=auto,
    xlabel style={at={(current axis.right of origin)},anchor=west},
    ylabel style={at={(current axis.above origin)},anchor=south},
    % extra x and y ticks for asymptotes
    extra x tick style={grid=major, grid style={dashed,black}},
    extra y tick style={grid=major, grid style={dashed,black}},
    /pgf/number format/1000 sep={},
  },    
  tick label style = {font=\footnotesize, fill=\currentbgcolor, inner sep=0pt},    
  xticklabel shift=1pt,    
  scaled ticks=false,
  scaled y ticks=false,
  unbounded coords=jump,
  samples=101,
  minor tick length=0,
  scatter/classes={
    open={mark=*,draw=black,fill=white, mark size=2, thin},
    closed={mark=*,draw=black,fill=black, mark size=2, thin}
  },
  width=3in,
}
\pgfplotsset{open/.style={only marks, mark=*, fill=white, thin, mark size=2}}
\pgfplotsset{closed/.style={only marks, mark=*, draw=black, fill=black,
    thin, mark size=2}}
\pgfplotsset{labeled/.style={nodes near coords, only marks, mark=*, draw=black, fill=black, thin, mark size=2, point meta=explicit symbolic}}

\usepackage{tikz}
\usetikzlibrary{
  matrix,%
  % enables the snake paths
  decorations.pathmorphing,%
  % enables bigger arrows
  decorations.markings,%
  decorations.pathreplacing,%
  decorations.text,%
  calc,%
  patterns,%
  shapes.geometric,%
  arrows,
  decorations.shapes,
  positioning,
  plotmarks,
  intersections
}
\tikzset{>=latex} % sets default arrow type in tikz
\tikzset{font=\normalsize}
\tikzset{mark size=1.5pt, mark options=thin}
\tikzset{pin distance=4pt,
  every pin edge/.style={<-, thin, shorten <= -2pt}}

\interfootnotelinepenalty=10000
\renewcommand{\thefootnote}{(\arabic{footnote})}

\usepackage{multirow, bigdelim}

% change to \usepackage[sols]{handout} to see solutions
\usepackage[]{handout}

\newcommand\iscurrentenumi[1]{%
  \ifthenelse{\equal{#1}{\theenumi}}{}{#1}}
\setenumerate[1]{ref=\#\arabic*}
\setenumerate[2]{ref=\protect\iscurrentenumi{\theenumi}(\alph*)}
\newcommand\enumiiprefix[2]{%
  \ifthenelse{{\equal{#1}{\theenumi}}\AND{\equal{#2}{(\alph{enumii})}}}{}{}%
  \ifthenelse{{\equal{#1}{\theenumi}}\AND{\NOT{\equal{#2}{(\alph{enumii})}}}}{#2}{}%
  \ifthenelse{\equal{#1}{\theenumi}}{}{#1#2}%
}
\setenumerate[3]{ref=\protect\enumiiprefix{\theenumi}{(\alph{enumii})}\roman*}
% Make an environment that puts items in columns, first going across. 
\usepackage{tabto}
\newenvironment{enumeratecols}[2][]
{\NumTabs{#2}%
  \begin{enumerate*}[%before={\hspace{\dimexpr-\parindent-1pt}\tab},%
    itemjoin={\tab},#1]}%
  {\end{enumerate*}}

\newlength{\tipmargin}
\makeatletter

\renewenvironment{framed}{
  \setlength{\tipmargin}{\textwidth-\linewidth}
  \advance\@totalleftmargin-\tipmargin
  \def\FrameCommand{\hspace{\tipmargin}\fbox}%
  \advance\linewidth\tipmargin
  \MakeFramed {\advance\hsize-\width \FrameRestore}%
}
{\endMakeFramed}

\makeatother

\definecolor{purple}{rgb}{0.5,0,1.0}

\usepackage{fancyhdr}
\lhead{\textcolor{white}{This content is protected and may not be shared, uploaded, or distributed.}}
\rhead{}
\rfoot{\vspace*{\fill}\footnotesize\textcopyright{} \emph{Harvard Math Ma}}
\renewcommand{\headrulewidth}{0pt}
\pagestyle{fancy}
\begin{document}

\htitle{Limits}

\begin{problemonly}
  \begin{framed}

You should be able to:
    \begin{itemize}[topsep=0pt,partopsep=0pt,parsep=8pt,itemsep=0pt]
    \item Calculate two-sided and one-sided limits of functions given algebraically or graphically, and understand why such a limit might not exist.

    \item Recognize when we need to factor and cancel to calculate a limit (for example, with $\displaystyle \lim_{x \to 2} \frac{x^2 - 4}{x - 2}$).

    \item Explain precisely what we mean when we write $\displaystyle \lim_{x \to a} f(x) = \infty$ or $\displaystyle \lim_{x \to a} f(x) = -\infty$. In particular, you should understand that, in both of these cases, $\displaystyle \lim_{x \to a} f(x)$ does not exist, but that saying $\displaystyle \lim_{x \to a} f(x) = \infty$ or $\displaystyle \lim_{x \to a} f(x) = -\infty$ gives more information than just saying ``$\displaystyle \lim_{x \to a} f(x)$ does not exist'' does.

    \item Use algebraic methods to determine whether a limit is $\pm \infty$.

    \end{itemize}
  \end{framed}
\end{problemonly}

\begin{enumerate}
\item  \label{limit-motivation}

  \begin{enumerate}
  \item

    \note{I'll start by having students remember the definition of the derivative, just to remember that it involves limits. Today, we'll finally get a better idea of what a limit is.}

    \begin{problem}[0.35in]
      If $g$ is a function, what's the definition of the derivative $g'(7)$?
    \end{problem}

    \begin{solution}
      $\displaystyle g'(7) = \lim_{h \to 0} \frac{g(7 + h) -
        g(7)}{h}$.
    \end{solution}

  \item

    \note{This is just to remind students how they computed such a limit. I will try to leave this up on the board so we can look back at it.}

    \begin{problem}[\vfill]
      If $g(x) = x^2$, calculate $g'(7)$. 
    \end{problem}

    \begin{solution}
      \begin{align*} 
        g'(7) &= \lim_{h \rightarrow 0} \frac{ g(7 +h) - g(7) } { 7+h - 7 } \\ 
        & = \lim_{h \rightarrow 0 } \frac{ (7+h)^2 -49} {h} \\ 
        & = \lim_{h \rightarrow 0} \frac{ 49 +14h+h^2 -49}{h} \\ 
        & = \lim_{h \rightarrow 0 } \frac{ 14h +h^2 } {h} \\ 
        & = \lim_{h \rightarrow 0} (14 +h)  \\
        &= 14
      \end{align*} 
    \end{solution}

  \end{enumerate}

\end{enumerate}
\note{I will probably say something about how we could also write $\displaystyle \lim_{w \to a} f(w)$ as $\displaystyle \lim_{\heartsuit \to a} f(\heartsuit)$.}

\begin{framed}
  \textbf{Limits.} You can think of the notation $\displaystyle
  \bm{\lim_{w \to a} f(w)}$ as shorthand for: ``What number (if any) does
  $f(w)$ approach as $w$ gets really close to $a$, without actually being
  equal to $a$?''  This number is called the \uline{limit of $f(w)$ as $w$
    approaches $a$}. 

  More precisely, saying $\displaystyle \lim_{w \to a} f(w) = L$ means
  that we can make the values of $f(w)$ as close as we like to $L$ by
  making $w$ sufficiently close to (but not equal to) $a$.
\end{framed}

\pspace{\newpage}

\begin{enumerate}[resume]
\item  \label{limit-graphical-examples-1-v2}

  \note{These are to emphasize how $\displaystyle \lim_{x \to a} f(x)$ is different from $f(a)$, as well as to introduce the desire for one-sided limits.}

  \begin{enumeratecols}{3}
  \item
    \begin{minipage}[t]{1.75in}\setlength{\parskip}{8pt}
      \raisebox{2ex-\height}{
        \begin{tikzpicture}
          \begin{axis}[xtick={-5, -3, ..., 5}, ytick={1, 2}, width=2.25in, height=1.6in]
            \addplot+[domain=-6:6, samples=200] {2*sin(deg(pi*(x-1)))/(pi*(x-1))} node[pos=0.6, right]{$f$};%
            \addplot+[open] coordinates {(1, 2)}; % 
          \end{axis}
        \end{tikzpicture}
      }

      \probonly{$\displaystyle \lim_{t \to 1} f(t) = \fitb{0.5in}{}$}
      \solonly{$\displaystyle \lim_{t \to 1} f(t) = \fitb{0.5in}{2}$}


      \probonly{$\displaystyle f(1) = \fitb{0.5in}{}$}
      \solonly{$\displaystyle f(1) = \fitb{0.5in}{\textbf{undefined}}$}
    \end{minipage}

  \item
    \begin{minipage}[t]{1.75in}\setlength{\parskip}{8pt}
      \raisebox{2ex-\height}{
        \begin{tikzpicture}
          \begin{axis}[xtick={-5, -3, ..., 5}, ytick={1, 2}, width=2.25in, height=1.6in]
            \addplot+[domain=-6:6, samples=200] {2*sin(deg(pi*(x-1)))/(pi*(x-1))} node[pos=0.6, right]{$f$};%
            \addplot+[open] coordinates {(1, 2)}; %
            \addplot+[closed] coordinates {(1, 1)}; % 
          \end{axis}
        \end{tikzpicture}
      }

      \probonly{$\displaystyle \lim_{x \to 1} f(x) = \fitb{0.5in}{}$}
      \solonly{$\displaystyle \lim_{x \to 1} f(x) = \fitb{0.5in}{2}$}
      

      \probonly{$\displaystyle f(1) = \fitb{0.5in}{}$}
      \solonly{$\displaystyle f(1) = \fitb{0.5in}{1}$}

      \probonly{$\displaystyle \lim_{x \to -3} f(x) = \fitb{0.5in}{}$}
      \solonly{$\displaystyle \lim_{x \to -3} f(x) = \fitb{0.5in}{0}$}

      \probonly{$\displaystyle f(-3) = \fitb{0.5in}{}$}
      \solonly{$\displaystyle f(-3) = \fitb{0.5in}{0}$}

    \end{minipage}

  \item

    \begin{minipage}[t]{1.85in}\setlength{\parskip}{8pt}
      \raisebox{2ex-\height}{
        \begin{tikzpicture}
          \begin{axis}[xtick={-6, -4, ..., 6}, ytick={-1, 1, 2}, width=2.25in, height=1.6in]
            \addplot+[domain=-6:-2.01, samples=200] {-sin(deg(pi*(x+2)))/(pi*(x+2))};%
            \addplot+[domain=-2:6] {0.5^x - 2} node[anchor=south east]{$f$}; %
            \addplot+[open] coordinates {(-2, 2)}; %
            \addplot+[closed] coordinates {(-2, -1)}; % 
          \end{axis}
        \end{tikzpicture}
      }

      \probonly{$\displaystyle \lim_{q \to -2} f(q) = \fitb{0.5in}{}$}
      \solonly{$\displaystyle \lim_{q \to -2} f(q) = \fitb{0.5in}{\textbf{DNE}}$}

      \probonly{$\displaystyle f(-2) = \fitb{0.5in}{}$}
      \solonly{$\displaystyle f(-2) = \fitb{0.5in}{-1}$}
    \end{minipage}

  \end{enumeratecols}

\end{enumerate}
\begin{framed}
  If $a$ is a number, $w$ can get really close to $a$ from two sides, the
  left and the right.  So, we also look at \uline{one-sided limits}:
  \begin{itemize}[topsep=0pt]

  \item The \uline{limit of $f(w)$ as $w$ approaches $a$ from the
      \emph{left}}, $\displaystyle \bm{\lim_{w \to a^-} f(w)}$, is the
    number that $f(w)$ approaches as $w$ gets really close to $a$, while
    remaining slightly \emph{less} than $a$.

  \item The \uline{limit of $f(w)$ as $w$ approaches $a$ from the
      \emph{right}}, $\displaystyle \bm{\lim_{w \to a^+} f(w)}$, is the
    number that $f(w)$ approaches as $w$ gets really close to $a$, while
    remaining slightly \emph{greater} than $a$.

  \end{itemize}
\end{framed}

\begin{enumerate}[resume]
\item 

  \begin{enumerate}
  \item

    \begin{problem}[\vfill]
      Sketch an example of a function $f$ for which $\displaystyle \lim_{x \to 2^-} f(x) = 1$ and $\displaystyle \lim_{x \to 2^+} f(x) = 3$. Can you sketch an example which also has the property that $\displaystyle \lim_{x \to 2} f(x)$ exists?
    \end{problem}

    \begin{solution}
      Here's one example, but there are many possibilities!
      \begin{center}
        \begin{tikzpicture}
          \begin{axis}[xtick={2}, ytick={1, 3} ]
            \addplot+[domain=2:4] {sqrt(x-1)+2}; %
            \addplot+[domain=-2:2] {x-1}; %
            \addplot[open] coordinates{(2,1)}; %
            \addplot[closed] coordinates{(2,3)}; %
          \end{axis}
        \end{tikzpicture}
      \end{center}
      It isn't possible to draw an example where $\displaystyle \lim_{x \to 2} f(x)$ exists.
    \end{solution}

  \item

    \note{This is to explicitly establish that, in order for $\displaystyle \lim_{x \to a} f(x)$ to exist, the one-sided limits must exist and equal each other.}

    \begin{problem}[\vfill\newpage]
      Generalize: In order for $\displaystyle \lim_{x \to a} f(x)$ to exist, what needs to happen with the one-sided limits $\displaystyle \lim_{x \to a^-} f(x)$ and $\displaystyle \lim_{x \to a^+} f(x)$?
    \end{problem}

    \begin{solution}
      $\displaystyle \lim_{x \to a^-} f(x)$ and $\displaystyle \lim_{x \to a^+} f(x)$ need to exist and be equal to each other.
    \end{solution}

  \end{enumerate}


\item 

  \note{Now, we'll think algebraically about some limits.

    In all of these limits, it makes sense to start by thinking about what happens to the numerator and denominator separately (and we can find out about these just by plugging in). In the first part, it's then straightforward to evaluate the limit. In the other two, we need to do more work.}

  Evaluate the following limits.

  \begin{enumerate}
  \item

    \begin{problem}[\stretch{0.5}]
      $\displaystyle \lim_{w \to 2} \frac{3w + 1}{w}$
    \end{problem}

    \begin{solution}
      We can use direct substitution to evaluate this limit.
      \[
        \lim_{w \to 2} \frac{3w+1}{w} = \frac{3(2)+1}{2}
        = \frac{7}{2}.
      \]
    \end{solution}

  \item

    \note{This is the same limit we evaluated in \ref{limit-motivation}.

      I don't expect students to realize that at first, and I'm interested to hear what they think a $\dfrac{0}{0}$ limit should be. I will remind them that what this is saying is that, when we plug in values of $h$ that are close to 0, we'll get $\dfrac{\text{something close to 0}}{\text{something close to 0}}$. In the Edfinity homework, they looked at expressions like $\dfrac{0.01}{0.00001}$ and $\dfrac{0.01}{0.01}$ and hopefully realized these have very different values. So, just knowing we have $\dfrac{\text{something close to 0}}{\text{something close to 0}}$ isn't enough information; we need to do more work.

      Then we'll look back at \ref{limit-motivation} and see that the method for evaluating such a limit is to factor and cancel.}

    \begin{problem}[\vfill]
      $\displaystyle \lim_{h \to 0} \dfrac{h^2 + 14h}{h}$
    \end{problem}

    \begin{solution}
      When we try to use direct substitution, we run into a problem. We get $\frac{0}{0}$, which is
      an undefined expression.
      \[
        \lim_{h \to 0} \frac{h^2+14h}{h} \stackrel?= \frac{0}{0}.
      \]
      You may have noticed that the same limit appeared in \ref{limit-motivation}. The strategy we used
      was to factor and cancel.
      \[
       \lim_{h \rightarrow 0 } \frac{ 14h +h^2 } {h} = \lim_{h \rightarrow 0} (14 +h) = 14.
      \]
    \end{solution}

  \item

    \begin{problem}[\vfill]
      $\displaystyle \lim_{q \to 2} \frac{2q - 4}{q^2 - 4}$
    \end{problem}

    \begin{solution}
      Again, we run into a problem when we try direct substitution.
      \[\lim_{q \to 2} \frac{2q-4}{q^2-4} \stackrel?= \frac{2(2)-4}{(2)^2-4} \stackrel?= \frac{0}{0}.\]
      The solution is again to factor and cancel.
      \begin{align*}
        \lim_{q \to 2} \frac{2q-4}{q^2-4} &= \lim_{q \to 2} \frac{2(q-2)}{(q-2)(q+2)}\\
        &= \lim_{q \to 2} \frac{2}{q+2} \\
        &= \frac{1}{2}.
      \end{align*}
    \end{solution}

  \end{enumerate}

\item  \label{limits-intro-to-infinity}

  \note{Introduction to infinite limits.}

  Sketch the graph of $f(x) = \frac{1}{x}$, and use it to find the following
  limits.

  \begin{enumerate}
  \item

    \note{It should be pretty intuitive to students to say this limit is $\infty$. Make sure students know that $\infty$ is not a number; when we say $\displaystyle \lim_{x \to 0^+} \tfrac{1}{x} = \infty$, we're actually saying, ``$\displaystyle \lim_{x \to 0^+} \tfrac{1}{x}$ does not exist because $\frac{1}{x}$ increases without bound as $x$ approaches 0 from the right.'' Students may wonder if they should write DNE or $\infty$; let them know we write $\infty$ because it's more informative.} 

    \begin{problem}
      $\displaystyle \lim_{x \to 0^+} \frac{1}{x}$
    \end{problem}

    \begin{solution}
      We know that the graph of $y=\frac{1}{x}$ has a vertical asymptote at $x=0$.
      
      \begin{tikzpicture}
        \begin{axis}[restrict y to domain=-5:5, xtick=\empty, ytick=\empty]
          \addplot[domain=-5:5] {1/x};
        \end{axis}
      \end{tikzpicture}

      We write $\lim\limits_{x \to 0^+} \frac{1}{x}= \infty$ as shorthand for
      ``the limit does not exist because the value of $\frac{1}{x}$ increases without bound as $x$
      approaches $0$ from the right.''
    \end{solution}

  \item
    \begin{problem}
      $\displaystyle \lim_{x \to 0^-} \frac{1}{x}$
    \end{problem}

    \begin{solution}
      We write $\lim\limits_{x \to 0^+} \frac{1}{x}= -\infty$ as shorthand for
      ``the limit does not exist because the value of $\frac{1}{x}$ decreases without bound as $x$
      approaches $0$ from the right.''
    \end{solution}

  \item
    \begin{problem}[\stretch{1.5}]
      $\displaystyle \lim_{x \to 0} \frac{1}{x}$
    \end{problem}

    \begin{solution}
      We say that $\lim\limits_{x \to 0} \frac{1}{x}$ does not exist.
    \end{solution}

  \end{enumerate}

  \note{Here, I want students to think along the lines of ``1 over a small positive number is a big positive number''. The main take-away is that we need to think about the sign of the denominator on either side of $x = 0$.} 

  \begin{problem}[\nstretch{1.5}]
    How could you have found these limits algebraically?
  \end{problem}

  \begin{solution}
    If we try to use direct substitution, we get an undefined expression: $\frac{1}{0}$. Unlike previous limits,
    factoring and canceling won't work. We should instead think about the value of $\frac{1}{x}$ as $x$ approaches
    $0$ from either side.

    When $x>0$, as $x$ takes on values closer and closer to zero,
    $\frac{1}{x}$ takes on higher and higher positive values,
    increasing without bound. This tells us that $\lim\limits_{x \to 0^+} \frac{1}{x}=\infty$.

    When $x<0$, as $x$ takes on values closer and closer to zero,
    $\frac{1}{x}$ takes on lower and lower negative values, decreasing without bound.
    This tells us that $\lim\limits_{x \to 0^-} \frac{1}{x} = -\infty$.

    Since $\frac{1}{x}$ increases without bound as $x$ approaches zero from the right and 
    decreases without bound as $x$ approaches zero from the left,
    $\lim\limits_{x \to 0} \frac{1}{x}$ does not exist.
    
  \end{solution}

\item 

  Find the following limits algebraically.

  \probonly{\begin{multicols}{3}}
    \begin{enumerate}
    \item
      \begin{problem}
        $\displaystyle \lim_{x \to 2^-} \frac{x - 5}{2 - x}$
      \end{problem}

      \begin{solution}
        When we attempt direct substitution, we get the undefined expression $\frac{-3}{0}$. 
        This is a good indication that the function is increasing or decreasing without bound
        as $x$ approaches $2$.
        
        When $x$ approaches $2$ from the left, the denominator $2-x$ takes on positive values closer
        and closer to zero.
        Since the numerator is negative, the fraction $\frac{x-5}{2-x}$ takes on lower
        and lower negative values, decreasing without bound. So $\lim\limits_{x \to 2^-} \frac{x-5}{2-x} = -\infty$.
      \end{solution}

    \item
      \begin{problem}
        $\displaystyle \lim_{x \to 2^+} \frac{x - 5}{2 - x}$
      \end{problem}

      \begin{solution}
        When $x$ approaches $2$ from the right, the denominator $2-x$ takes on negative values closer and
        closer to zero.
        Since the numerator is also negative, the fraction $\frac{x-5}{2-x}$ takes on higher and higher
        positive values, increasing without bound. So $\lim\limits_{x \to 2^-} \frac{x-5}{2-x} = \infty$.
      \end{solution}

    \item      
      \begin{problem}
        $\displaystyle \lim_{x \to 2} \frac{x - 5}{2 - x}$
      \end{problem}

      \begin{solution}
        When $x$ approaches $2$, the function increases without bound on one side and decreases without
        bound on the other, so $\lim\limits_{x \to 2} \frac{x-5}{2-x}$ does not exist.
      \end{solution}

    \end{enumerate}
  \probonly{\end{multicols}}

  \pspace{\stretch{1.5}}

\item 
  Find the following limits algebraically. If the limit does not exist, can you say anything about the one-sided limits?

  \begin{enumerate}
  \item
    \begin{problem}[\vfill]
      $\displaystyle \lim_{t \to 3} \frac{t - 4}{t - 3}$
    \end{problem}

    \begin{solution}
      When we attempt direct substitution, we get the undefined expression $\frac{-1}{0}$. This is a good
      indication that the function is increasing or decreasing without bound as $t$ approaches $3$.
      
      
      As $t$ approaches $3$ from either side, the numerator is negative.
      When $t$ approaches $3$ from the left, the denominator is also negative, so 
      $\lim\limits_{t \to 0^-} \frac{t-4}{t-3} = \infty$. When $t$ approaches $3$ from the right, the denominator
      is positive, so $\lim\limits_{t \to 0^-} \frac{t-4}{t-3} = -\infty$. 
      This means $\lim\limits_{t \to 0} \frac{t-4}{t-3}$ does not exist.
    \end{solution}

  \item
    \begin{problem}[\vfill]
      $\displaystyle \lim_{t \to 4} \frac{t - 4}{t - 3}$
    \end{problem}

    \begin{solution}
      When we attempt direct substitution, we get the expression $\frac{0}{1}=0$, which is not undefined.
      Thus, $\lim\limits_{t \to 4} \frac{t-4}{t-3}=0$. 
    \end{solution}

  \item
    \begin{problem}[\vfill]
      $\displaystyle \lim_{x \to 1} \frac{x^2 - 6x + 8}{x - 2}$
    \end{problem}

    \begin{solution}
      Direct substitution gives us $\frac{3}{-1}=-3$, so 
      $\lim\limits_{t \to 1}{(x^2-6x+8)}{x-2}=-3$. 
    \end{solution}

  \item
    \begin{problem}[\vfill\newpage]
      $\displaystyle \lim_{x \to 2} \frac{x^2 - 6x + 8}{x - 2}$
    \end{problem}

    \begin{solution}
      Direct substitution gives us $\frac{0}{0}$, which prompts us to try to factor and cancel.
      \begin{align*}
          \lim_{x \to 2} \frac{x^2-6x+8}{x-2} &= \lim_{x \to 2} \frac{(x-2)(x-4)}{x-2} \\
          &= \lim_{x \to 2} (x-4) \\
          &= -2.
      \end{align*}
    \end{solution}

  \end{enumerate}

\item  \label{graphical-limits-1}

  {
    The graph of a function $g(x)$ is given below.  

    \newcommand{\graphcommand}[1]{
      \begin{tikzpicture}
        \begin{axis}[xmin=-4, xmax=4, xtick={-4, -3, ..., 4}, ytick={-3,
            ..., 3}, enlargelimits=true, ymin=-2, ymax=2, width=3.5in,
          height=2.25in, axis equal=true]
          \draw[thick] (-4, 0) arc[radius=1, start angle=270, end
          angle=360]; %
          \draw[thick] (-3, 1) -- (-1, 2); %
          \addplot+[sharp plot] coordinates { (-1, 1) (1, -2) (2, 0) (4,
            2) }; %
          \addplot+[open] coordinates { (-1, 1) (1, -2) (3, 1) }; %
          \addplot+[closed] coordinates { (-4, 0) (-1, 2) (1, -1) (3, 2)
            (4, 2)}; %
          #1 %
        \end{axis}
      \end{tikzpicture}
    }

    \graphcommand{}

  }

  Evaluate the following, or explain why they do not exist. 
  \probonly{\begin{multicols}{3}}
    \begin{enumerate}
    \item
      \begin{problem}
        $\displaystyle \lim_{x \rightarrow 1^+} g(x)$
      \end{problem}

      \begin{solution}
          $\displaystyle \lim_{x \rightarrow 1^+} g(x) = -2$
      \end{solution}


    \item
      \begin{problem}
        $\displaystyle \lim_{x \rightarrow 1^-} g(x)$
      \end{problem}

      \begin{solution}
          $\displaystyle \lim_{x \rightarrow 1^-} g(x) = -2$
      \end{solution}

    \item
      \begin{problem}
        $\displaystyle \lim_{x \rightarrow 1} g(x)$
      \end{problem}

      \begin{solution}
          $\displaystyle \lim_{x \rightarrow 1} g(x) = -2$
      \end{solution}

    \item
      \begin{problem}
        $g(1)$
      \end{problem}

      \begin{solution}
        $g(1) = -1$
      \end{solution}

    \item
      \begin{problem}
        $\displaystyle \lim_{x \rightarrow -4^+} g(x)$
      \end{problem}

      \begin{solution}
          $\displaystyle \lim_{x \rightarrow -4^+} g(x)=0$
      \end{solution}
      

    \item
      \begin{problem}
        $\displaystyle \lim_{x \rightarrow -3} g(x)$
      \end{problem}

      \begin{solution}
        $\displaystyle \lim_{x \rightarrow -3} g(x)=1$
      \end{solution}

    \item
      \begin{problem}
        $\displaystyle \lim_{x \rightarrow -1} g(x)$
      \end{problem}

      \begin{solution}
        $\displaystyle \lim_{x \rightarrow -1} g(x)$ does not exist because the one-sided limits
        are not equal.
      \end{solution}

    \item
      \begin{problem}
        $g(-1)$
      \end{problem}

      \begin{solution}
       $g(-1)=2$
      \end{solution}

    \item
      \begin{problem}
        $\displaystyle \lim_{x \rightarrow 3} g(x)$
      \end{problem}

      \begin{solution}
        $\displaystyle \lim_{x \rightarrow 3} g(x) = 1$
      \end{solution}

    \item
      \begin{problem}
        $g(3)$
      \end{problem}

      \begin{solution}
        $g(3) = 2$
      \end{solution}

    \item
      \begin{problem}
        $\displaystyle \lim_{x \rightarrow 4^-} g(x)$
      \end{problem}

      \begin{solution}
        $\displaystyle \lim_{x \rightarrow 4^-} g(x)=2$
      \end{solution}

    \end{enumerate}
  \probonly{\end{multicols}}

  \pspace{1in}

\item 

  \note{This is to establish two things:
    \begin{itemize}
    \item When we look at a limit as $x \to a$, we only care what happens when $x$ is extremely close to $a$.
    \item Once we've decided which is the right piece to consider, we can just plug $x = a$ in to evaluate the limit. 
    \end{itemize}
  }

  \begin{problem}[\vfill]
    Define
    $\displaystyle f(x) =
      \begin{cases} 
        x^2 & \text{ if } x \leq 3.9999 \\
        \frac{1}{2} x + 3 & \text{ if } 3.9999 < x < 4.0001 \\
        \frac{1}{x} & \text{ if } x \geq 4.0001    
      \end{cases}$

    What is $\displaystyle \lim_{x \to 4} f(x)$?
  \end{problem}

  \begin{solution}
    To evaluate $\lim\limits_{x \to 4} f(x)$, we need to consider what happens as $x$ approaches $4$.
    Since $f$ is piecewise defined, we should decide which pieces are relevant. In this case,
    the only piece which is relevant is the one for the interval $3.9999 < x < 4.0001$. If we
    zoom in on the function at $x=4$, we would eventually see only this piece.

    Now that we've identified the relevant piece, we can use direct substitution to evaluate the limit:
    \begin{align*}
      \lim_{x \to 4} f(x) &= \lim_{x \to 4} \tfrac{1}{2} x + 3 \\
      &= \tfrac{1}{2}(4) + 3 \\
      &= 5.
    \end{align*}
  \end{solution}
\end{enumerate}
\end{document}
